\section{AMC 12 12 2000}
        \begin{enumerate}
        	\item (\textit{AMC 12 12 2000}) In the year $2001$, the United States will host the International Mathematical Olympiad. Let $I,M,$ and $O$ be distinct positive integers such that the product $I \cdot M \cdot O = 2001$. What is the largest possible value of the sum $I + M + O$?


 $\textbf{(A)}\ 23 \qquad \textbf{(B)}\ 55 \qquad \textbf{(C)}\ 99 \qquad \textbf{(D)}\ 111 \qquad \textbf{(E)}\ 671$


	\item (\textit{AMC 12 12 2000}) $2000(2000^{2000}) =$
$\textbf{(A)}\ 2000^{2001} \qquad \textbf{(B)}\ 4000^{2000} \qquad \textbf{(C)}\ 2000^{4000} \qquad \textbf{(D)}\ 4,000,000^{2000} \qquad \textbf{(E)}\ 2000^{4,000,000}$


	\item (\textit{AMC 12 12 2000}) Each day, Jenny ate $20\%$ of the jellybeans that were in her jar at the beginning of that day. At the end of the second day, $32$ remained. How many jellybeans were in the jar originally?


 $\textbf{(A)}\ 40 \qquad \textbf{(B)}\ 50 \qquad \textbf{(C)}\ 55 \qquad \textbf{(D)}\ 60 \qquad \textbf{(E)}\ 75$


	\item (\textit{AMC 12 12 2000}) The Fibonacci sequence $1,1,2,3,5,8,13,21,\ldots$ starts with two $1$'s, and each term afterwards is the sum of its two predecessors. Which one of the ten digits is the last to appear in the units position of a number in the Fibonacci sequence?


 $\textbf{(A)}\ 0 \qquad \textbf{(B)}\ 4 \qquad \textbf{(C)}\ 6 \qquad \textbf{(D)}\ 7 \qquad \textbf{(E)}\ 9$


	\item (\textit{AMC 12 12 2000}) If $|x - 2| = p,$ where $x < 2,$ then $x - p =$
$\textbf{(A)}\ -2 \qquad \textbf{(B)}\ 2 \qquad \textbf{(C)}\ 2-2p \qquad \textbf{(D)}\ 2p-2 \qquad \textbf{(E)}\ |2p-2|$


	\item (\textit{AMC 12 12 2000}) Two different prime numbers between $4$ and $18$ are chosen. When their sum is subtracted from their product, which of the following numbers could be obtained?


 $\textbf{(A)}\ 22 \qquad \textbf{(B)}\ 60 \qquad \textbf{(C)}\ 119 \qquad \textbf{(D)}\ 194 \qquad \textbf{(E)}\ 231$


	\item (\textit{AMC 12 12 2000}) How many positive integers $b$ have the property that $\log_{b} 729$ is a positive integer?


 $\textbf{(A)}\ 0 \qquad \textbf{(B)}\ 1 \qquad \textbf{(C)}\ 2 \qquad \textbf{(D)}\ 3 \qquad \textbf{(E)}\ 4$


	\item (\textit{AMC 12 12 2000}) Figures $0$, $1$, $2$, and $3$ consist of $1$, $5$, $13$, and $25$ nonoverlapping unit squares, respectively. If the pattern were continued, how many nonoverlapping unit squares would there be in figure 100?


 
                    \begin{figure}[H]
                        \centering
                        \includegraphics[width=0.5\textwidth]{Problem/AMC_12_12_2000_8_Problem_Figure_1.png}
                    \end{figure}
                    
$\textbf{(A)}\ 10401 \qquad\textbf{(B)}\ 19801 \qquad\textbf{(C)}\ 20201 \qquad\textbf{(D)}\ 39801 \qquad\textbf{(E)}\ 40801$


	\item (\textit{AMC 12 12 2000}) Mrs. Walter gave an exam in a mathematics class of five students. She entered the scores in random order into a spreadsheet, which recalculated the class average after each score was entered. Mrs. Walter noticed that after each score was entered, the average was always an integer. The scores (listed in ascending order) were $71,76,80,82,$ and $91$. What was the last score Mrs. Walters entered?


 $\textbf{(A)} \ 71 \qquad \textbf{(B)} \ 76 \qquad \textbf{(C)} \ 80 \qquad \textbf{(D)} \ 82 \qquad \textbf{(E)} \ 91$


	\item (\textit{AMC 12 12 2000}) The point $P = (1,2,3)$ is reflected in the $xy$-plane, then its image $Q$ is rotated $180^\circ$ about the $x$-axis to produce $R$, and finally, $R$ is translated  $5$ units in the positive-$y$ direction to produce $S$. What are the coordinates of $S$?


 $\textbf {(A) } (1,7, - 3) \qquad \textbf {(B) } ( - 1,7, - 3) \qquad \textbf {(C) } ( - 1, - 2,8) \qquad \textbf {(D) } ( - 1,3,3) \qquad \textbf {(E) } (1,3,3)$


	\item (\textit{AMC 12 12 2000}) Two non-zero real numbers, $a$ and $b,$ satisfy $ab = a - b$. Which of the following is a possible value of $\frac {a}{b} + \frac {b}{a} - ab$?


 $\textbf{(A)} \ - 2 \qquad \textbf{(B)} \ \frac {- 1}{2} \qquad \textbf{(C)} \ \frac {1}{3} \qquad \textbf{(D)} \ \frac {1}{2} \qquad \textbf{(E)} \ 2$


	\item (\textit{AMC 12 12 2000}) Let $A, M,$ and $C$ be nonnegative integers such that $A + M + C=12$. What is the maximum value of $A \cdot M \cdot C + A \cdot M + M \cdot C + A \cdot C$?


 $\textbf{(A)}\ 62 \qquad \textbf{(B)}\ 72 \qquad \textbf{(C)}\ 92 \qquad \textbf{(D)}\ 102 \qquad \textbf{(E)}\ 112$


	\item (\textit{AMC 12 12 2000}) One morning each member of Angela’s family drank an $8$-ounce mixture of coffee with milk. The amounts of coffee and milk varied from cup to cup, but were never zero. Angela drank a quarter of the total amount of milk and a sixth of the total amount of coffee. How many people are in the family?


 $\textbf {(A)}\ 3 \qquad \textbf {(B)}\ 4 \qquad \textbf {(C)}\ 5 \qquad \textbf {(D)}\ 6 \qquad \textbf {(E)}\ 7$


	\item (\textit{AMC 12 12 2000}) When the \href{https://artofproblemsolving.com/wiki/index.php?title=Mean}{mean}, \href{https://artofproblemsolving.com/wiki/index.php?title=Median}{median}, and \href{https://artofproblemsolving.com/wiki/index.php?title=Mode}{mode} of the list


 
 \[10,2,5,2,4,2,x\]


 are arranged in increasing order, they form a non-constant \href{https://artofproblemsolving.com/wiki/index.php?title=Arithmetic\_progression}{arithmetic progression}. What is the sum of all possible real values of $x$?


 $\textbf {(A)}\ 3 \qquad \textbf {(B)}\ 6 \qquad \textbf{(C)}\ 9 \qquad \textbf {(D)}\ 17 \qquad \textbf {(E)}\ 20$


	\item (\textit{AMC 12 12 2000}) Let $f$ be a \href{https://artofproblemsolving.com/wiki/index.php?title=Function}{function} for which $f(x/3) = x^2 + x + 1$. Find the sum of all values of $z$ for which $f(3z) = 7$.


 $\textbf {(A)}\ -1/3 \qquad \textbf {(B)}\ -1/9 \qquad \textbf {(C)}\ 0 \qquad \textbf {(D)}\ 5/9 \qquad \textbf {(E)}\ 5/3$


	\item (\textit{AMC 12 12 2000}) A checkerboard of $13$ rows and $17$ columns has a number written in each square, beginning in the upper left corner, so that the first row is numbered $1,2,\ldots,17$, the second row $18,19,\ldots,34$, and so on down the board. If the board is renumbered so that the left column, top to bottom, is $1,2,\ldots,13$, the second column $14,15,\ldots,26$ and so on across the board, some squares have the same numbers in both numbering systems. Find the sum of the numbers in these squares (under either system). 


 $\textbf {(A)}\ 222 \qquad \textbf {(B)}\ 333\qquad \textbf {(C)}\ 444 \qquad \textbf {(D)}\ 555 \qquad \textbf {(E)}\ 666$


	\item (\textit{AMC 12 12 2000}) 
                    \begin{figure}[H]
                        \centering
                        \includegraphics[width=0.5\textwidth]{Problem/AMC_12_12_2000_17_Problem_Figure_1.png}
                    \end{figure}
                    A \href{https://artofproblemsolving.com/wiki/index.php?title=Circle}{circle} centered at $O$ has \href{https://artofproblemsolving.com/wiki/index.php?title=Radius}{radius} $1$ and contains the point $A$. The segment $AB$ is \href{https://artofproblemsolving.com/wiki/index.php?title=Tangent\_(geometry)}{tangent} to the circle at $A$ and $\angle AOB = \theta$. If point $C$ lies on $\overline{OA}$ and $\overline{BC}$ bisects $\angle ABO$, then $OC =$
$\textbf {(A)}\ \sec^2 \theta - \tan \theta \qquad \textbf {(B)}\ \frac 12 \qquad \textbf {(C)}\ \frac{\cos^2 \theta}{1 + \sin \theta}\qquad \textbf {(D)}\ \frac{1}{1+\sin\theta} \qquad \textbf {(E)}\ \frac{\sin \theta}{\cos^2 \theta}$


	\item (\textit{AMC 12 12 2000}) In year $N$, the $300$th day of the year is a Tuesday. In year $N+1$, the $200$th day is also a Tuesday. On what day of the week did the $100$th day of year $N-1$ occur?


 $\textbf {(A)}\ \text{Thursday} \qquad \textbf {(B)}\ \text{Friday}\qquad \textbf {(C)}\ \text{Saturday}\qquad \textbf {(D)}\ \text{Sunday}\qquad \textbf {(E)}\ \text{Monday}$


	\item (\textit{AMC 12 12 2000}) In \href{https://artofproblemsolving.com/wiki/index.php?title=Triangle}{triangle} $ABC$, $AB = 13$, $BC = 14$, $AC = 15$. Let $D$ denote the \href{https://artofproblemsolving.com/wiki/index.php?title=Midpoint}{midpoint} of $\overline{BC}$ and let $E$ denote the intersection of $\overline{BC}$ with the \href{https://artofproblemsolving.com/wiki/index.php?title=Angle\_bisector}{bisector} of angle $BAC$. Which of the following is closest to the area of the triangle $ADE$?


 $\textbf {(A)}\ 2 \qquad \textbf {(B)}\ 2.5 \qquad \textbf {(C)}\ 3 \qquad \textbf {(D)}\ 3.5 \qquad \textbf {(E)}\ 4$


	\item (\textit{AMC 12 12 2000}) If $x,y,$ and $z$ are positive numbers satisfying $x + \frac{1}{y} = 4, y + \frac{1}{z} = 1,$ and $z + \frac{1}{x} = \frac73,$ then what is the value of $xyz$ ?


 $\textbf {(A)}\ 2/3 \qquad \textbf {(B)}\ 1 \qquad \textbf {(C)}\ 4/3 \qquad \textbf {(D)}\ 2 \qquad \textbf {(E)}\ 7/3$


	\item (\textit{AMC 12 12 2000}) Through a point on the \href{https://artofproblemsolving.com/wiki/index.php?title=Hypotenuse}{hypotenuse} of a \href{https://artofproblemsolving.com/wiki/index.php?title=Right\_triangle}{right triangle}, lines are drawn \href{https://artofproblemsolving.com/wiki/index.php?title=Parallel}{parallel} to the legs of the triangle so that the triangle is divided into a \href{https://artofproblemsolving.com/wiki/index.php?title=Square}{square} and two smaller right triangles. The area of one of the two small right triangles is $m$ times the area of the square. The \href{https://artofproblemsolving.com/wiki/index.php?title=Ratio}{ratio} of the area of the other small right triangle to the area of the square is 


 $\textbf {(A)}\ \frac{1}{2m+1} \qquad \textbf {(B)}\ m \qquad \textbf {(C)}\ 1-m \qquad \textbf {(D)}\ \frac{1}{4m} \qquad \textbf {(E)}\ \frac{1}{8m^2}$


	\item (\textit{AMC 12 12 2000}) The graph below shows a portion of the curve defined by the quartic polynomial $P(x) = x^4 + ax^3 + bx^2 + cx + d$. Which of the following is the smallest?


 $\textbf{(A)}\ P(-1)\\ \textbf{(B)}\ \text{The\ product\ of\ the\ zeros\ of\ } P\\ \textbf{(C)}\ \text{The\ product\ of\ the\ non-real\ zeros\ of\ } P \\ \textbf{(D)}\ \text{The\ sum\ of\ the\ coefficients\ of\ } P \\ \textbf{(E)}\ \text{The\ sum\ of\ the\ real\ zeros\ of\ } P$

                    \begin{figure}[H]
                        \centering
                        \includegraphics[width=0.5\textwidth]{Problem/AMC_12_12_2000_22_Problem_Figure_1.png}
                    \end{figure}
                    

	\item (\textit{AMC 12 12 2000}) Professor Gamble buys a lottery ticket, which requires that he pick six different integers from $1$  through $46$, inclusive. He chooses his numbers so that the sum of the base-ten logarithms of his six numbers is an integer. It so happens that the integers on the winning ticket have the same property— the sum of the base-ten logarithms is an integer. What is the probability that Professor Gamble holds the winning ticket?


 $\textbf {(A)}\ 1/5 \qquad \textbf {(B)}\ 1/4 \qquad \textbf {(C)}\ 1/3 \qquad \textbf {(D)}\ 1/2 \qquad \textbf {(E)}\ 1$


	\item (\textit{AMC 12 12 2000}) 
                    \begin{figure}[H]
                        \centering
                        \includegraphics[width=0.5\textwidth]{Problem/AMC_12_12_2000_24_Problem_Figure_1.png}
                    \end{figure}
                    If circular \href{https://artofproblemsolving.com/wiki/index.php?title=Arc}{arcs} $AC$ and $BC$ have \href{https://artofproblemsolving.com/wiki/index.php?title=Center}{centers} at $B$ and $A$, respectively, then there exists a \href{https://artofproblemsolving.com/wiki/index.php?title=Circle}{circle} \href{https://artofproblemsolving.com/wiki/index.php?title=Tangent\_(geometry)}{tangent} to both $\stackrel{\frown}{AC}$ and $\stackrel{\frown}{BC}$, and to $\overline{AB}$. If the length of $\stackrel{\frown}{BC}$ is $12$, then the \href{https://artofproblemsolving.com/wiki/index.php?title=Circumference}{circumference} of the circle is 


 $\textbf {(A)}\ 24 \qquad \textbf {(B)}\ 25 \qquad \textbf {(C)}\ 26 \qquad \textbf {(D)}\ 27 \qquad \textbf {(E)}\ 28$


	\item (\textit{AMC 12 12 2000}) Eight congruent \href{https://artofproblemsolving.com/wiki/index.php?title=Equilateral\_triangle}{equilateral triangles}, each of a different color, are used to construct a regular \href{https://artofproblemsolving.com/wiki/index.php?title=Octahedron}{octahedron}. How many distinguishable ways are there to construct the octahedron? (Two colored octahedrons are distinguishable if neither can be rotated to look just like the other.) 


 $\textbf {(A)}\ 210 \qquad \textbf {(B)}\ 560 \qquad \textbf {(C)}\ 840 \qquad \textbf {(D)}\ 1260 \qquad \textbf {(E)}\ 1680$

                    \begin{figure}[H]
                        \centering
                        \includegraphics[width=0.5\textwidth]{Problem/AMC_12_12_2000_25_Problem_Figure_1.png}
                    \end{figure}
                    


\end{enumerate}